% Fichier complété par tools/complete_missing_corrections.py.
% Source : SLCI/SLCI-07-Ordre12/542_Cours/542_Cours.tex
% Statut : rédigé (3 question(s)).
\begin{corrigebox}[Corrigé rédigé]
\CorrigeQuestion{1}
$H(p)=\dfrac{K\omega_0^2}{p^2+2\xi\omega_0p+\omega_0^2}
            =\dfrac{K}{1+2\xi p/\omega_0+p^2/\omega_0^2}$.
\CorrigeQuestion{2}
$\xi>1$ : réponse apériodique sans dépassement; $\xi=1$ : régime critique;
            $0<\xi<1$ : oscillations amorties; $\xi=0$ : oscillations non amorties.
            Un dépassement apparaît pour $\xi<1$ et augmente lorsque $\xi$ diminue.
\CorrigeQuestion{3}
Le gain est lu sur la valeur finale. Pour un régime pseudo-périodique,
            $D_1=e^{-\pi\xi/\sqrt{1-\xi^2}}$ donne $\xi$, et la pseudo-période
            $T_p=2\pi/(\omega_0\sqrt{1-\xi^2})$ donne $\omega_0$.
            Sans dépassement, on ajuste deux pôles réels à partir des temps caractéristiques.
\end{corrigebox}
